Tato kapitola poskytuje krátký, praktický úvod do části „AI v přípravě výuky“ s cílem nabídnout okamžitě použitelné kroky a odkazy pro garanty předmětů a vyučující na VUT. Najdete zde popis klíčových kompetencí, postupů pro rychlý pilot v kurzu a konkrétní inspiraci k nástrojům, které lze vyzkoušet ve výuce.

\medskip

\textbf{Cíle kapitoly.} Po přečtení této kapitoly by měl čtenář být schopen:
\begin{itemize}
  \item identifikovat oblasti předmětu vhodné pro doplnění AI‑podpory (příprava materiálů, cvičení, automatická zpětná vazba),
  \item použít připravené šablony a prompty k rychlé tvorbě sylabu, 90minutové lekce a pracovních listů (viz Přílohy),
  \item navrhnout jednoduchý pilotní workflow včetně základních kroků pro GDPR/etiku a transparentní komunikaci se studenty (viz Kapitola 2).
\end{itemize}
Poznámka: položky kompetencí jsou stručným souhrnem praktických cílů kapitoly (general background knowledge).

\medskip

\textbf{Jak kapitolu využít při přípravě předmětu — doporučený postup (rychlý checklist).}
\begin{enumerate}
  \item Definujte výukové výstupy, u kterých očekáváte přidanou hodnotu od AI (např. více variant cvičení, rychlejší formativní zpětná vazba).
  \item Vyberte jednu konkrétní oblast (např. generování pracovních listů nebo tvorba kvízů) a zvolte nástroj / šablonu z Příloh pro pilot.
  \item Připravte malý pilot: jeden týden / jedna lekce / malá skupina studentů; předem definujte metriky úspěchu (čas úspory, spokojenost studentů, kvalita materiálů).
  \item Ověřte požadavky na ochranu osobních údajů a transparentnost (informujte studenty, jaký nástroj se používá a jak deklarovat použití AI) — viz Kapitola 2 a šablony do sylabu.
  \item Po pilotu proveďte rychlou reflexi a upravte šablony/prompty; pokud pilot uspěje, rozšiřte nasazení postupně.
\end{enumerate}
(Body 3–5 jsou praktické doporučení odvozené z implementačních vzorů v příručce; konkretizace opatření k ochraně dat je rozvedena v Kapitole 2.)

\medskip

\textbf{Krátká inspirace — konkrétní nástroje a příklady k vyzkoušení.}
\begin{itemize}
  \item Pro názorné uvedení principů strojového učení a AI ve výuce lze využít vzdělávací svět Minecraft Education; v praxi jsou dostupné připravené hodiny (Hour of Code), které pomáhají žákům pochopit základní principy AI a programování \cite{kb_ai_101_tipu_pdf_04e4a115_page_15_2}.
  \item Při práci s matematikou a procvičováním postupů doporučujeme vyzkoušet funkce typu „Pokrok v matematice“, které generují příklady a mohou vyžadovat nahrání postupu řešení — užitečné pro hodnocení procesu, nikoli jen výsledku \cite{kb_ai_101_tipu_pdf_04e4a115_page_8_1}.
  \item Pokud chcete sledovat technické novinky a oficiální learning‑materiály k nástrojům (např. Microsoft Copilot, Azure, další), doporučeným zdrojem je portál Microsoft Learn jako systematický průvodce a dokumentace \cite{kb_ai_101_tipu_pdf_04e4a115_page_54_1}.
\end{itemize}

\medskip

\textbf{Krátké pokyny k etice a GDPR (shrnutí).} V této kapitole pouze navádíme, kde etiku a ochranu dat zapracovat do pracovního postupu: vždy zahrňte před pilotem základní informaci studentům o použití AI, zvažte nutnost DPIA pro zpracování osobních údajů a použijte šablony pro DPA/DPIA uvedené v Přílohách. Podrobná doporučení a checklisty jsou k dispozici v Kapitole 2 a v části Přílohy (Vzor DPIA, AI Policy Template).

\medskip

\textbf{Kde najdete šablony a další materiály.} Kopírovatelné prompty, šablony pro sylaby a checklisty najdete v části Přílohy (Prompt Library, AI Policy Template, Vzor DPIA). Online repozitář s video tutoriály a mechanismem kvartálních aktualizací je uveden v závěru příručky (viz informace o repozitáři a aktualizacích).